\section{Results}
\label{sec:results}

% Every number in this section is generated by `imp sweep` / `imp harness` +
% `imp tables` from committed run configs; tables are \input from
% paper/tables/. Do not hand-enter values.

Throughout, \emph{the full estimator} names \populacefit{} with all three
design choices on (weighted, gated, chained); each table indents its
ablations and variants beneath it. The \microimpute{} forest, OLS, quantile
regression, and the hot decks are the standard-practice baselines
(Table~\ref{tab:methods}), and the weighted marginal draw is the
unconditional lower bound. A row appears only where its method applies: the
net-worth-derived variants at production scale, the constrained hot deck on
the within-survey tasks, the random hot deck wherever the donor weights are
not uniform.
Unless a caption says otherwise, every cell is the mean over the ten paired
seeds with the standard deviation across seeds in parentheses, and bold
marks the best non-reference value in a column: lowest for distances and
losses, highest for coverage, closest to 0.5 for the two-sample AUC,
closest to 1 for tail ratios, and closest to the held-out truth's own value
for fragility.

\subsection{Wealth within the SCF}
\label{sec:results-scf}

\begin{table}[t]
\centering
\footnotesize
\resizebox{\textwidth}{!}{% Generated by `imp tables` -- do not edit by hand. SCF wealth task, mean (sd) over 10 paired seeds.
\begin{tabular}{lrrrrr}
\hline
Method & Debt pinball & Debt $W_1$ & Debt zero-share err. & Net worth pinball & Net worth $W_1$ \\
\hline
microcosm-fit & 49,201 (2,637) & 13,001 (6,285) & 0.020 (0.012) & 551,317 (36,012) & \textbf{135,871} (51,106) \\
\quad -- unweighted & 50,218 (2,452) & 32,247 (16,930) & 0.012 (0.014) & 703,457 (77,173) & 829,158 (445,108) \\
\quad -- unchained & 49,201 (2,637) & 13,001 (6,285) & 0.020 (0.012) & 556,622 (23,349) & 158,606 (44,310) \\
\quad -- ungated & 49,149 (2,832) & 12,714 (4,817) & 0.072 (0.029) & 563,223 (35,315) & 180,233 (52,533) \\
\quad -- ungated, unchained & 49,149 (2,832) & 12,714 (4,817) & 0.072 (0.029) & 571,907 (35,101) & 180,570 (86,541) \\
QRF (microimpute) & 50,433 (3,120) & 16,619 (5,358) & 0.018 (0.018) & 560,819 (36,903) & 234,367 (142,057) \\
OLS & 70,053 (17,637) & 156,295 (74,822) & 0.226 (0.018) & 1,116,572 (383,007) & 3,565,923 (2,062,063) \\
Quantile regression & 49,797 (3,391) & 29,460 (20,456) & 0.151 (0.031) & 534,420 (37,801) & 295,470 (110,986) \\
NND hot deck (py-statmatch) & 49,537 (2,627) & 13,770 (5,141) & 0.017 (0.013) & 644,921 (51,255) & 509,013 (399,757) \\
Weighted random hot deck & 48,802 (2,595) & \textbf{11,395} (4,789) & 0.020 (0.014) & 506,558 (29,427) & 172,984 (48,701) \\
Constrained hot deck & 49,518 (2,706) & 16,244 (6,243) & \textbf{0.011} (0.007) & 642,310 (44,183) & 642,051 (529,395) \\
Weighted marginal draw & \textbf{48,277} (2,586) & 14,617 (4,732) & 0.020 (0.017) & \textbf{443,554} (20,651) & 215,398 (87,264) \\
\hline
\end{tabular}
}
\caption{Wealth imputation within the \scf{} (debt and net worth from
demographics and income): weighted pinball loss, Wasserstein-1 to the weighted
donor distribution, and zero-share error, mean (sd) over ten paired seeds.
Bold marks the best non-reference value per column. Lower is better.}
\label{tab:scf-wealth}
\end{table}

Table~\ref{tab:scf-wealth} reports the within-\scf{} task. The headline is the
weighting ablation: switching the full estimator to \texttt{weights="none"}
degrades net-worth Wasserstein-1 by a factor of six (135,871 to 829,158) and
debt by a factor of 2.5 (13,001 to 32,247). The \scf{}'s list-sample design
oversamples wealthy households at low weight, so an unweighted fit learns the
sample's wealth distribution rather than the population's --- this is the
weight-blindness failure mode of Section~1 measured directly, and no other
design choice comes close to its effect size. Removing the gates raises
net-worth $W_1$ by a third (to 180,233) and removing chaining by 17 percent
(to 158,606); both paired deltas are smaller than their own standard
deviation across seeds (Table~\ref{tab:ablations}). What the gates change
here is the zero mass: with them off, debt zero-share error rises from 2.0
to 7.2 percentage points.

Against matching, the comparison turns on how the donor is chosen. The
nearest-neighbour hot deck nearly ties the full estimator on debt marginal
fit (13,770 vs.\ 13,001) but is 3.7 times worse on net
worth (509,013). The weighted random hot deck, which draws among the nearest
donors in proportion to their weights, is not: it posts the lowest debt
$W_1$ of any method (11,395) and 172,984 on net worth, 1.3 times the full
estimator and inside its seed-to-seed spread. The weights do that, and not
the draw: the random hot deck, which makes the same draw with every near
donor equally likely, closes none of the gap (1,033,450 on net worth).
The constrained hot deck, whose cap counts receivers and ignores the
weights, does no better than the nearest-neighbour one (642,051). Within
the survey, then, the full estimator's large margin over nearest-neighbour
matching is a margin over weight-blind donor choice, not over matching: the
one hot deck whose choice of donor uses the weights closes most of it, and
the three whose choice ignores them do not.
Among the other baselines,
the \microimpute{} forest sits at 234,367 on net worth, and OLS is
catastrophic on both targets (156,295 and 3.6 million): normal residuals cannot
represent a wealth distribution. Linear quantile regression leaves 15
percent of the debt zero mass unaccounted for, against 7 percent for the
gates-off ablations and at most 2 percent for the other forest-based
methods.

One result reads wrong until the metric is examined: the \emph{unconditional}
weighted marginal draw posts the best pinball loss on both targets. Pinball
loss as computed here is a marginal score whose predicted quantiles are read
without receiver weights (Section~\ref{sec:metrics}). In a survey that
oversamples the wealthy, a conditional method's pooled draw follows the
receiver sample's composition, while the marginal draw is centered on the
weighted population by construction, so this column favors the marginal
draw for a reason unrelated to conditional fit, and we do not read its lead
as evidence about the methods. More generally, no receiver-side marginal
metric can distinguish a method that draws from the correct weighted
marginal while ignoring predictors from one that also gets the conditionals
right. That is what the population-view harness of
Section~\ref{sec:results-harness}, which scores the joint, is for.

\subsection{Zero-inflated components and reweighting fragility}
\label{sec:results-cps}

\begin{table}[t]
\centering
\footnotesize
\resizebox{\textwidth}{!}{% Generated by `imp tables` -- do not edit by hand. CPS zero-inflated components, mean (sd) over 10 paired seeds.
\begin{tabular}{lrrrr}
\hline
Method & Interest pinball & Interest zero-share & Dividend pinball & Dividend zero-share \\
\hline
microcosm-fit & 1,491 (63.58) & 0.013 (0.011) & 320.68 (45.26) & 0.006 (0.004) \\
\quad -- unweighted & 1,491 (63.62) & 0.011 (0.007) & 320.68 (45.29) & 0.008 (0.005) \\
\quad -- unchained & 1,491 (63.58) & 0.013 (0.011) & 320.69 (45.26) & \textbf{0.006} (0.005) \\
\quad -- ungated & 1,491 (63.68) & 0.024 (0.015) & 321.05 (45.27) & 0.072 (0.006) \\
\quad -- ungated, unchained & 1,491 (63.68) & 0.024 (0.015) & 320.99 (45.26) & 0.043 (0.009) \\
QRF (microimpute) & 1,493 (63.34) & 0.014 (0.011) & 322.17 (45.24) & 0.044 (0.010) \\
OLS & 2,527 (63.57) & 0.422 (0.012) & 961.52 (112.29) & 0.844 (0.005) \\
Quantile regression & 1,493 (63.23) & 0.270 (0.016) & 320.91 (45.19) & 0.251 (0.011) \\
NND hot deck (py-statmatch) & 1,493 (63.89) & 0.021 (0.019) & 320.94 (45.22) & 0.009 (0.009) \\
Weighted random hot deck & \textbf{1,491} (63.49) & \textbf{0.009} (0.009) & 320.68 (45.28) & 0.007 (0.008) \\
Constrained hot deck & 1,491 (63.38) & 0.018 (0.013) & 320.75 (45.27) & 0.007 (0.005) \\
Weighted marginal draw & 1,491 (63.61) & 0.011 (0.013) & \textbf{320.68} (45.25) & 0.006 (0.004) \\
\hline
\end{tabular}
}
\caption{Zero-inflated income components in the raw \cps{} (interest and
dividend income for adults, on a fixed subsample of 20,000 of the file's
108,044 adults): weighted pinball loss, Wasserstein-1 to the weighted donor
distribution, and zero-share error, mean
(sd) over ten paired seeds. Bold marks the best non-reference value per
column; lower is better.}
\label{tab:cps-components}
\end{table}

\begin{table}[t]
\centering
\footnotesize
\resizebox{\textwidth}{!}{% Generated by `imp tables` -- do not edit by hand. Worst-case single-record aggregate share at kappa=5.
\begin{tabular}{lrrrr}
\hline
Method & Interest fragility & (truth) & Dividend fragility & (truth) \\
\hline
microcosm-fit & 0.384 (0.051) & 0.372 (0.043) & 0.675 (0.100) & 0.614 (0.104) \\
\quad -- unweighted & 0.395 (0.043) & 0.372 (0.043) & 0.699 (0.057) & 0.614 (0.104) \\
\quad -- unchained & 0.384 (0.051) & 0.372 (0.043) & 0.661 (0.100) & 0.614 (0.104) \\
\quad -- ungated & 0.384 (0.043) & 0.372 (0.043) & 0.781 (0.052) & 0.614 (0.104) \\
\quad -- ungated, unchained & 0.384 (0.043) & 0.372 (0.043) & 0.716 (0.068) & 0.614 (0.104) \\
QRF (microimpute) & 0.364 (0.064) & 0.372 (0.043) & 0.633 (0.087) & 0.614 (0.104) \\
OLS & 0.055 (0.010) & 0.372 (0.043) & 0.055 (0.009) & 0.614 (0.104) \\
Quantile regression & \textbf{0.368} (0.105) & 0.372 (0.043) & 0.444 (0.073) & 0.614 (0.104) \\
NND hot deck (py-statmatch) & 0.391 (0.027) & 0.372 (0.043) & 0.647 (0.140) & 0.614 (0.104) \\
Random hot deck & 0.431 (0.120) & 0.372 (0.043) & 0.736 (0.119) & 0.614 (0.104) \\
Weighted random hot deck & 0.410 (0.085) & 0.372 (0.043) & 0.761 (0.066) & 0.614 (0.104) \\
Constrained hot deck & 0.379 (0.044) & 0.372 (0.043) & \textbf{0.602} (0.131) & 0.614 (0.104) \\
Weighted marginal draw & 0.356 (0.042) & 0.372 (0.043) & 0.746 (0.067) & 0.614 (0.104) \\
\hline
\end{tabular}
}
\caption{Reweighting fragility at $\kappa=5$: worst-case single-record share
of each component's aggregate over bounded multiplicative reweightings, next
to the held-out truth's own fragility. The target is the truth's exposure ---
bold marks the closest method; substantially lower means the method fails to
generate realistic tails, substantially higher means it plants landmines.}
\label{tab:fragility}
\end{table}

Table~\ref{tab:cps-components} reports the \cps{} components task. Pinball
loss barely separates the non-OLS methods --- as a marginal score it is not
built to --- but the zero mass separates
them sharply. The full estimator holds zero-share error to 0.6 and 1.3
percentage points (dividends, interest). With its gates off, the same
estimator misses by 7.2 and 2.4 points, and its dividend $W_1$ doubles (223
against 102). It errs by drawing too many zeros: in a forest grown to
purity an atom at zero forms larger leaves than the positive values do, and
a forest that pools leaf samples across trees over-represents the atom in
each row's conditional. The gate replaces that with a classifier's
probability. The \microimpute{} forest, which has no gate either and also
differs in its sampler (Section~\ref{sec:ablations}), misses the dividend
zero share by 4.4 points; linear quantile regression misses by 25--27 points because a
99-level linear quantile grid cannot hold a five-sixths zero mass; and OLS
misses the dividend zero share by 84 points --- with normal residual draws,
nearly every adult becomes a dividend recipient. The hot decks, which donate
observed values, hold the zero mass within 0.7--2.1 points, as expected. On
marginal fit the nearest-neighbour hot deck's $W_1$ is more than twice the
full estimator's on both components (227 and 760 against 102 and 290),
while the other three hot decks match or beat it within seed noise (random
125 and 299; weighted random 107 and 232; constrained 98 and 231). The
pattern differs from the \scf{}'s in one respect. There the hot decks whose
choice ignores the weights did no better than the nearest-neighbour one;
here they do as well as the weighted one.
With person weights that carry little outcome information, what the
nearest-donor rule loses on this task it loses to donor reuse, which the
cap and the random draw both limit, and not to weight-blindness.

Table~\ref{tab:fragility} reports the landmine diagnostic. Two readings
matter. First, the held-out truth is itself concentrated: under $\kappa=5$
reweighting a single record can carry 37 percent of the interest aggregate
and 61 percent of the dividend aggregate at this sample size --- heavy-tailed
capital income is genuinely fragile, and a faithful method should
\emph{match} that exposure, not minimize it. Second, every forest- and
donation-based method tracks the truth's fragility closely, while OLS sits
far below it, at 0.06 against 0.37 and 0.61 --- understating the true
exposure by factors of roughly 7 and 11. Its aggregate is spread across many
moderate records because the method never generates realistic tails; the
zero-share and Wasserstein numbers show the same failure from other angles.
The highest dividend exposures belong to the gates-off ablation (0.78), the
weighted random hot deck (0.76), and the marginal draw (0.75), against a
truth of 0.61 whose own standard deviation across seeds is 0.10; none
approaches a landmine on this within-survey task.
The landmine pathology of Section~1 arises when unweighted cross-file fits
meet downstream reweighting, and the direct evidence for it here is the tail
block of the population-view experiment
(Section~\ref{sec:results-harness}).

\subsection{Ablations: attributing the gains}
\label{sec:results-ablations}

\begin{table}[t]
\centering
\footnotesize
\resizebox{\textwidth}{!}{% Generated by `imp tables` -- do not edit by hand. Ablation minus candidate, paired-seed delta as mean (sd); positive = ablation worse (lower-is-better metrics).
\begin{tabular}{lllrrrr}
\hline
Task & Target & Metric & $\Delta$ unweighted & $\Delta$ unchained & $\Delta$ ungated & $\Delta$ ungated, unchained \\
\hline
cps components & dividend\_income & pinball\_loss & 0.000 (0.056) & 0.005 (0.033) & 0.361 (0.097) & 0.302 (0.089) \\
cps components & dividend\_income & wasserstein1 & 36.11 (46.53) & 11.29 (94.58) & 120.57 (55.98) & 107.34 (74.80) \\
cps components & dividend\_income & zero\_share\_error & 0.002 (0.006) & -0.000 (0.006) & 0.067 (0.008) & 0.037 (0.010) \\
cps components & interest\_income & pinball\_loss & -0.016 (0.307) & 0.000 (0.000) & -0.024 (0.347) & -0.024 (0.347) \\
cps components & interest\_income & wasserstein1 & -26.28 (114.17) & 0.000 (0.000) & -11.50 (187.77) & -11.50 (187.77) \\
cps components & interest\_income & zero\_share\_error & -0.002 (0.006) & 0.000 (0.000) & 0.011 (0.023) & 0.011 (0.023) \\
scf wealth & debt & pinball\_loss & 1,017 (709.88) & 0.000 (0.000) & -51.77 (595.03) & -51.77 (595.03) \\
scf wealth & debt & wasserstein1 & 19,246 (16,086) & 0.000 (0.000) & -287.07 (7,733) & -287.07 (7,733) \\
scf wealth & debt & zero\_share\_error & -0.008 (0.011) & 0.000 (0.000) & 0.052 (0.034) & 0.052 (0.034) \\
scf wealth & networth & pinball\_loss & 152,140 (51,569) & 5,305 (23,838) & 11,906 (23,086) & 20,590 (19,445) \\
scf wealth & networth & wasserstein1 & 693,287 (431,527) & 22,735 (70,684) & 44,362 (67,012) & 44,699 (93,246) \\
scf wealth & networth & zero\_share\_error & -0.000 (0.001) & 0.000 (0.001) & 0.001 (0.003) & 0.001 (0.002) \\
\hline
\end{tabular}
}
\caption{Paired-seed deltas from the full estimator, mean (sd) over ten
seeds, for its four ablations: weights off, chaining off, gates off, and
gates and chaining off. Positive means the ablation is worse on these
lower-is-better metrics; a delta near zero means the knocked-out design
choice is inert on that task. On each task's first target (interest, debt)
there is nothing to chain on, so the chaining column is zero and the two
gates-off columns agree.}
\label{tab:ablations}
\end{table}

Table~\ref{tab:ablations} attributes the full estimator's behavior to its
design choices. Each cell answers one question: on the same ten splits, how
much did this metric worsen when one design choice was knocked out? The
$+693{,}287$ under $\Delta$ unweighted on net-worth $W_1$, for instance,
reads: refitting the identical estimator without weights raised the
marginal-fit distance by \$693,287 on average. Because ablation and
estimator share each split, the delta subtracts out split-to-split noise;
the parenthetical standard deviation says how much the delta itself varies
across the ten splits.

The ordering is task-dependent in an informative way. On the \scf{} ---
informative weights, heavy tails --- weighting dominates ($+693{,}287$
net-worth $W_1$ when removed). Gating ($+44{,}362$, sd $67{,}012$) and
chaining ($+22{,}735$, sd $70{,}684$) shift net-worth $W_1$ by less than
their dispersion across seeds, and gating's clear effect is on the debt
zero share ($+0.052$, sd $0.034$). On the \cps{} components --- weakly
informative weights, extreme zero inflation --- weighting has a small effect
(removing it raises dividend $W_1$ by 36, sd 47, from 102, and leaves the
other metrics within seed noise), chaining is inert, and gating carries the
load: with the gates off, dividend zero-share error rises by 6.7 percentage
points (sd 0.8) and dividend $W_1$ by 121 (sd 56). Interest, whose zero share
is nearer one half, moves less (1.1 points, sd 2.3). Which choice binds depends on the survey's
design and the target's structure. Where a choice does not bind, its cost is
small: the unweighted ablation is slightly ahead on \scf{} debt zero share
(0.012 against 0.020, Table~\ref{tab:scf-wealth}) and on the uniform-weight
OpenML suite (Section~\ref{sec:results-openml}).

\subsection{Cross-dataset transfer}
\label{sec:results-openml}

\begin{table}[t]
\centering
\footnotesize
\resizebox{\textwidth}{!}{% Generated by `imp tables` -- do not edit by hand. Wasserstein-1 to donor, mean over 10 paired seeds; lower is better.
\begin{tabular}{lrrrrrrr}
\hline
Method & abalone & brazilian\_houses & elevators & house\_sales & onlinenewspopularity & space\_ga & Mean rank \\
\hline
microcosm-fit & 0.254 & 166.64 & 0.000 & 7,080 & 274.67 & 0.008 & 4.42 \\
\quad -- unweighted & 0.257 & 184.92 & 0.000 & 6,744 & 233.10 & 0.008 & 4.00 \\
\quad -- ungated & 0.254 & 166.64 & 0.000 & 7,080 & 274.67 & 0.009 & 4.58 \\
QRF (microimpute) & 0.174 & 191.36 & 0.000 & 8,461 & 546.16 & 0.012 & 5.67 \\
OLS & 0.532 & 212.09 & 0.002 & 90,518 & 6,117 & 0.015 & 9.83 \\
Quantile regression & 0.321 & 212.04 & 0.001 & 24,115 & 382.98 & 0.085 & 9.00 \\
NND hot deck (py-statmatch) & 0.176 & 199.12 & 0.000 & 10,212 & 161.43 & 0.010 & 4.67 \\
Weighted random hot deck & 0.237 & 196.93 & 0.000 & 16,570 & 217.64 & 0.010 & 5.33 \\
Constrained hot deck & 0.151 & 198.10 & 0.000 & 9,248 & 144.90 & 0.009 & \textbf{3.50} \\
Weighted marginal draw & 0.148 & 210.50 & 0.000 & 7,135 & 143.52 & 0.011 & 4.00 \\
\hline
\end{tabular}
}
\caption{Wasserstein-1 to the donor distribution across the six OpenML
regression datasets (uniform weights), mean over ten paired seeds, with mean
rank across datasets among the ten methods shown. The two unchained arms are
omitted: with a single target they coincide with their chained
counterparts. So is the random hot deck, which under uniform weights is the
weighted random one. The two datasets with
more than 20,000 rows (\texttt{house\_sales}, 21,613, and
\texttt{onlinenewspopularity}, 39,644) are capped at 20,000 by a fixed
uniform subsample.}
\label{tab:openml}
\end{table}

Table~\ref{tab:openml} runs the surface over the six OpenML regression
datasets under uniform weights. Two observations. First, with weights
uninformative by construction, the full estimator and its unweighted ablation
differ only by a uniform resample of the training rows, and they rank in the
upper half (mean ranks 4.4 and 4.0) alongside the unconditional marginal draw
(4.0); the linear methods rank last (9.0 and 9.8). The gates-off ablation
equals the full estimator on five of the six datasets, whose targets have a
single sign class and so no gate to remove. Second --- and more important for
this paper's argument --- a method with no conditional structure at all ties
for second of the ten on a \emph{marginal} distance, ahead of the full
estimator, and the best mean rank (3.5) goes to the constrained hot deck,
which here uses each donor at most once and so reproduces the donor marginal
almost by construction. Marginal metrics saturate: they reward drawing
from the right marginal and cannot rank what imputation is actually for.
The population-view harness exists because of exactly this ceiling.

\subsection{The SCF\texorpdfstring{$\to$}{->}CPS population view}
\label{sec:results-harness}

\begin{table}[t]
\centering
\footnotesize
\resizebox{\textwidth}{!}{% Generated by `imp tables` -- do not edit by hand. SCF->CPS population view vs held-out SCF, mean (sd) over 10 donor splits. AUC of 0.5 is indistinguishable; q99 ratio of 1 is a perfect tail. Minimal profile: 6 shared predictors, 2 targets, ASEC 2025 receiver.
\begin{tabular}{lrrrrr}
\hline
Method & Energy distance & Coverage & C2ST AUC & NW $W_1$/sd & NW q99 ratio \\
\hline
SCF sample (floor) & 0.009 (0.006) & 0.987 (0.009) & 0.496 (0.023) & 0.013 (0.003) & 1.10 (0.202) \\
\hline
microcosm-fit & \textbf{0.139} (0.026) & 0.960 (0.014) & 0.732 (0.013) & 0.031 (0.007) & 0.696 (0.112) \\
\quad -- unweighted & 0.140 (0.026) & 0.958 (0.014) & \textbf{0.731} (0.016) & 0.159 (0.041) & 1.93 (0.255) \\
\quad -- unchained & 0.139 (0.026) & 0.958 (0.016) & 0.742 (0.013) & 0.029 (0.009) & 0.765 (0.177) \\
QRF (microimpute) & 0.140 (0.027) & 0.957 (0.012) & 0.775 (0.008) & 0.029 (0.007) & 0.805 (0.156) \\
OLS & 0.166 (0.045) & 0.732 (0.087) & 0.988 (0.005) & 0.434 (0.282) & \textbf{0.967} (0.505) \\
Quantile regression & 0.141 (0.027) & 0.938 (0.028) & 0.927 (0.018) & 0.052 (0.013) & 0.530 (0.118) \\
NND hot deck (py-statmatch) & 0.141 (0.027) & \textbf{0.962} (0.013) & 0.753 (0.014) & 0.051 (0.016) & 0.762 (0.148) \\
Weighted marginal draw & 0.139 (0.026) & 0.937 (0.005) & 0.848 (0.008) & \textbf{0.014} (0.004) & 1.09 (0.222) \\
\hline
\end{tabular}
}
\caption{The population-view experiment, \emph{minimal} profile (six shared
predictors, two targets, a fixed subsample of 20,000 of the 55,762 ASEC 2025
receiver households): candidates assembled from real
\cps{} demographics plus each method's imputed wealth, under \cps{}
household weights, scored against the held-out \scf{} through the \scf{}
view. The first row is the sampling-noise floor; cells are mean (sd) over
the ten paired donor splits. Energy distance lower is better; coverage
higher; C2ST AUC of 0.5 is indistinguishable; net-worth $W_1$/sd lower is
better; a q99 ratio of 1 is a perfect extreme tail. Bold marks the best
non-reference value per column under those conventions.}
\label{tab:harness}
\end{table}

\begin{table}[t]
\centering
\footnotesize
\resizebox{\textwidth}{!}{% Generated by `imp tables` -- do not edit by hand. SCF->CPS population view vs held-out SCF, mean (sd) over 10 donor splits. AUC of 0.5 is indistinguishable; q99 ratio of 1 is a perfect tail. Production-scale profile: 10 shared predictors, 4 chained targets, pooled ASEC 2023-2025 receiver.
\begin{tabular}{lrrrrr}
\hline
Method & Energy distance & Coverage & C2ST AUC & NW $W_1$/sd & NW q99 ratio \\
\hline
SCF sample (floor) & 0.012 (0.005) & 0.986 (0.003) & 0.499 (0.021) & 0.013 (0.003) & 1.10 (0.202) \\
\hline
microcosm-fit & 0.118 (0.024) & 0.969 (0.009) & 0.761 (0.026) & 0.027 (0.008) & 0.767 (0.190) \\
\quad -- unweighted & 0.120 (0.024) & \textbf{0.971} (0.008) & 0.754 (0.011) & 0.133 (0.038) & 2.25 (0.443) \\
\quad -- unchained & 0.118 (0.024) & 0.967 (0.011) & 0.819 (0.015) & 0.028 (0.006) & 0.699 (0.112) \\
\quad -- ungated & \textbf{0.117} (0.024) & 0.970 (0.011) & 0.758 (0.013) & 0.028 (0.007) & 0.767 (0.157) \\
\quad -- ungated, unchained & 0.117 (0.024) & 0.967 (0.012) & 0.831 (0.012) & 0.026 (0.008) & 0.778 (0.129) \\
\quad -- net worth derived & 0.118 (0.024) & 0.970 (0.009) & \textbf{0.744} (0.011) & 0.027 (0.008) & 0.773 (0.186) \\
\quad -- unchained, net worth derived & 0.118 (0.024) & 0.970 (0.011) & 0.756 (0.014) & 0.031 (0.008) & 0.764 (0.184) \\
QRF (microimpute) & 0.126 (0.025) & 0.965 (0.011) & 0.817 (0.009) & 0.033 (0.022) & 1.15 (0.445) \\
OLS & 0.152 (0.059) & 0.877 (0.055) & 0.995 (0.002) & 0.435 (0.282) & \textbf{0.963} (0.503) \\
Quantile regression & 0.120 (0.026) & 0.951 (0.011) & 0.970 (0.007) & 0.051 (0.013) & 0.485 (0.098) \\
NND hot deck (py-statmatch) & 0.119 (0.025) & 0.971 (0.010) & 0.768 (0.009) & 0.024 (0.008) & 0.747 (0.107) \\
Random hot deck & 0.121 (0.025) & 0.968 (0.010) & 0.782 (0.008) & 0.073 (0.020) & 1.57 (0.182) \\
Weighted random hot deck & 0.118 (0.025) & 0.968 (0.011) & 0.777 (0.006) & 0.037 (0.010) & 0.598 (0.110) \\
Weighted marginal draw & 0.119 (0.024) & 0.931 (0.015) & 0.968 (0.003) & \textbf{0.015} (0.003) & 1.12 (0.241) \\
\hline
\end{tabular}
}
\caption{The population-view experiment, \emph{production-scale} profile (ten
shared predictors, four chain-ordered balance-sheet targets, a fixed
subsample of 20,000 of the 168,852 pooled ASEC 2023--2025 receiver
households). Columns, cell format, and bolding as in
Table~\ref{tab:harness}.}
\label{tab:harness-scale}
\end{table}

Tables \ref{tab:harness} and \ref{tab:harness-scale} report the
population-view experiment in its two profiles. The floor behaves as designed
in both: another sample of the same survey is statistically indistinguishable
(energy 0.009 and 0.012, coverage 0.987 and 0.986, AUC 0.496 and 0.499), and
its q99 ratio of 1.10 calibrates how much a 99th-percentile ratio moves under
sampling noise alone. Four findings follow.

First, \emph{the weighting failure is a tail failure, and only the tail
block sees it}. The unweighted ablation ties the full estimator on energy
distance, coverage, and the C2ST in both profiles --- and its imputed
net-worth q99 is 1.93 times the holdout's in the minimal profile and 2.25
times at production scale, with tail $W_1$/sd five times the full
estimator's (0.159 vs.\ 0.031 minimal; 0.133 vs.\ 0.027 at scale). Fitting the \scf{}
unweighted learns the wealthy oversample's conditionals, and the resulting
inflation concentrates beyond the resolution of capped pairwise geometry: we
report this as a finding about evaluation practice. Any harness whose joint
metrics subsample --- which at these sample sizes they must --- needs an
uncapped tail block for heavy-tailed economic variables, or it will certify
files whose top percentile is wrong by a factor of two.

Second, \emph{chaining's contribution appears only where the targets are
jointly constrained, only a joint metric sees it, and most of it is the
constraint itself}. In the
minimal profile (two targets) the unchained ablation is indistinguishable
from the full estimator (AUC 0.742 against 0.732, under half the floor's
seed-to-seed standard deviation). At production scale, where the four targets satisfy
the accounting identity networth = fin + nfin $-$ debt exactly in the donor,
fitting them independently costs 5.8 points of C2ST AUC (0.819 vs.\ 0.761)
while leaving every marginal and tail statistic essentially unchanged ---
the classifier detects balance sheets that do not add up, which no marginal
metric can. The net-worth-derived variants show how much of that is the
identity. Imputing the three components independently and computing net
worth from them scores 0.756, level with the fully chained estimator.
Chaining the three components as well scores 0.744, the lowest AUC of any
method. Chaining therefore adds 1.2 AUC points beyond what the identity
supplies, and deriving net worth is 1.7 points better than imputing it as a
fourth chained target. Gating does not register on the joint scores of
either profile: every gates-off row sits within one floor standard
deviation of its gated counterpart (at scale, AUC 0.758 chained and 0.831
unchained, against 0.761 and 0.819).

Third, \emph{richer conditioning makes the classifier two-sample test ---
the scorecard's omnibus block --- decisive}. The
unconditional weighted marginal draw --- at or near the top on marginal
metrics throughout this paper --- moves from AUC 0.848 in the
minimal profile to 0.968 at production scale, while the full estimator stays
within two points of the lowest AUC in each profile (0.732 against its
unweighted ablation's 0.731; 0.761 against its net-worth-derived variant's
0.744). With ten shared predictors and a
four-component wealth block, a classifier finds the missing
demographics--wealth coupling almost perfectly. This is the empirical case
for carrying strong auxiliary predictors in a production imputation stack:
richer conditioning is what turns joint-structure failures from undetectable
into obvious.

Fourth, \emph{hot-deck matching is a near-peer on this task, and drawing
among near donors does not improve on the nearest one here}. The
nearest-neighbour hot deck sits within noise of the full estimator on energy and coverage in both
profiles, within 2.1 and 0.7 AUC points, and its donated values slightly
beat the full estimator's on tail $W_1$ at scale (0.024 vs.\ 0.027). The
weighted random hot deck, which closed the gap within the \scf{}, does not
carry that over: its AUC is higher than the nearest-neighbour hot deck's in
both profiles (0.761 and 0.777, against 0.753 and 0.768), and its tail fit
is worse at scale (0.037). The random hot deck separates the draw from the
weights again. Its AUC is about the weighted one's (0.764 and 0.782), so
the higher AUC comes with the draw. Its q99 ratio is about twice the
weighted one's (1.08 and 1.57, against 0.52 and 0.60), so the weights lower
the imputed tail, as they do for the estimator in the first finding. What
the full estimator adds over matching is
therefore not within-survey marginal fit, once the donor is chosen with the
weights (Table~\ref{tab:scf-wealth}), and not the zero mass, which donation
preserves (Table~\ref{tab:cps-components}). It is a modest margin on the
joint. In the minimal profile its AUC is 2.1 points below the
nearest-neighbour hot deck's and 2.9 below the weighted random one's (0.732
against 0.753 and 0.761). At scale it ties the
nearest-neighbour hot deck as specified, and sits 2.4 points below it once
net worth is derived (0.744 against 0.768), so that both files' balance
sheets add up. Each of these margins is of the order of one floor standard
deviation (0.023 and 0.021).

One more tail observation cuts against the full
estimator: its imputed net-worth q99 falls short of the holdout's, at a
ratio of 0.70 in the minimal profile and 0.77 at production scale. The
nearest-neighbour hot deck, which donates observed values, shows the same
shortfall (0.76 and 0.75), and the weighted random hot deck falls shorter
still (0.52 and 0.60), while the unconditional marginal draw, which draws
from the donor's weighted marginal without conditioning, sits near 1 (1.09
and 1.12, against a floor of 1.10). The shortfall is therefore shared by a
forest and two donation methods that all condition on the receiver's
demographics, and this experiment does not isolate its cause. The random
hot deck's higher ratios (1.08 and 1.57) come from leaving the weights out
of its draw, and sit on top of the same shortfall. No method in
this surface gets the extreme tail right for free; the tail block is what
keeps that visible.

The arms that distinguish these findings --- the gates-off ablations, the
net-worth-derived variants, and the random and weighted random hot decks
--- were also run on the 2025 wave of the \scf{}
(Tables~\ref{tab:scf-wealth-2025}--\ref{tab:harness-scale-2025}), outside
the pre-registered replication of Appendix~\ref{app:scf2025}. Each pattern
repeats there: the debt zero share drifts with the gates off, the weighted
random hot deck closes most of the nearest-neighbour hot deck's
within-survey gap on net worth where the random one leaves most of it open,
and at production scale the identity
recovers the unchained loss while chaining adds about a point more.
